\exposetitle{I}{Algebraic structures. Cohomology of groups}
\label{I}

\begin{center}
\textit{by M.~Demazure}
\end{center}

{\renewcommand{\thefootnote}{(0)}\nde{Version of 13/10/2024}}%
This expos\'e consists of two parts; the first collects a certain number
of general definitions and sets down notations which will often be used
in the sequel, the second treats the cohomology of groups and leads to
Theorem~5.3.3 (vanishing of the cohomology of diagonalizable groups).

We choose once and for all a
Universe.%
{\renewcommand{\thefootnote}{(1)}\nde{\Cf SGA~4, Exp.~I, \S0 and the
Appendix; see also the discussion in \textup{[DG70]}, p.~xxvi.}}
All the definitions introduced and all the constructions carried out
will be relative to this Universe. We shall systematically allow
ourselves the following abuse of language: in order to define a functor
\(f\colon\mathcal{C}\to\mathcal{C}'\), we shall content ourselves with
defining the object \(f(S)\) of \(\mathcal{C}'\) for every object \(S\)
of \(\mathcal{C}\), whenever there will be no ambiguity as to the manner
of defining \(f(h)\) for an arrow \(h\) of \(\mathcal{C}\). In practice,
we shall say: let \(f\colon\mathcal{C}\to\mathcal{C}'\) be the functor
defined by \(f(S)=\cdots\).

\sgasection{1}{Generalities}
\label{I.sec.1}

\numpara{1.1}
Let \(\mathcal{C}\) be a category. One will denote by
\(\widehat{\mathcal{C}}\) the category
\(\Hom(\mathcal{C}^{\circ},\Ens)\) of contravariant functors from
\(\mathcal{C}\) to the category \(\Ens\) of
sets.%
{\renewcommand{\thefootnote}{(2)}\nde{One calls it the category of
presheaves on \(\mathcal{C}\), \cf IV~4.3.1.}}
There exists a canonical functor
\(\mathbf{h}\colon\mathcal{C}\to\widehat{\mathcal{C}}\) which associates
with every \(X\in\Ob(\mathcal{C})\) the functor \(\mathbf{h}_{X}\) such
that
\[
\mathbf{h}_{X}(S)=\Hom(S,X).
\]
For every functor \(\mathbf{F}\in\Ob(\widehat{\mathcal{C}})\), one defines
(\cf for example EGA \(0_{\mathrm{III}}\), 8.1.4) a
bijection%
{\renewcommand{\thefootnote}{(3)}\nde{This result is often called the
``Yoneda lemma''; we shall use this terminology in other N.D.E.}}
\[
\Hom(\mathbf{h}_{X},\mathbf{F})\isomto\mathbf{F}(X).
\]

In particular, for every pair \(X\), \(Y\) of objects of \(\mathcal{C}\),
the canonical map below is bijective:
\[
\Hom_{\mathcal{C}}(X,Y)\isomto
\Hom_{\widehat{\mathcal{C}}}(\mathbf{h}_{X},\mathbf{h}_{Y});
\]
\ie the functor \(\mathbf{h}\) is \emph{fully faithful}. It therefore
defines an isomorphism of \(\mathcal{C}\) onto a \emph{full} subcategory of
\(\widehat{\mathcal{C}}\), and an equivalence of \(\mathcal{C}\) with the
\emph{full} subcategory of \(\widehat{\mathcal{C}}\) formed of the
\emph{representable} functors (\ie isomorphic to a functor of the form
\(\mathbf{h}_{X}\)). In the sequel, we shall often identify \(X\) and
\(\mathbf{h}_{X}\). The following numbers are intended to show that this
identification can be made without danger.

\begin{remark}{1.1.1}
\label{I.1.1.1}
{\renewcommand{\thefootnote}{(4)}\nde{This remark has been added.}}%
One sometimes needs the following variant. Let \(\mathcal{D}\) be a
\emph{full} subcategory of \(\mathcal{C}\) and let
\(X,Y\in\Ob(\mathcal{D})\); denote by \(\mathbf{h}'_{X}\) and
\(\mathbf{h}'_{Y}\) the restrictions to \(\mathcal{D}\) of
\(\mathbf{h}_{X}\) and \(\mathbf{h}_{Y}\). Then one has
\[
\Hom_{\mathcal{C}}(X,Y)=\Hom_{\mathcal{D}}(X,Y)
=\Hom_{\widehat{\mathcal{D}}}(\mathbf{h}'_{X},\mathbf{h}'_{Y})
\]
and hence: to give a morphism \(X\to Y\) ``is the same thing'' as to
give, in a manner functorial in \(T\), a map
\(\varphi(T)\colon\Hom(T,X)\to\Hom(T,Y)\), for every
\(T\in\Ob(\mathcal{D})\).
\end{remark}

\numpara{1.2}
{\renewcommand{\thefootnote}{(5)}\nde{The order has been modified, so as
to introduce fibered products before monomorphisms, \cf N.D.E.~(9).}}%
One will say that \(\mathbf{F}\) is a \emph{subobject} (or a
\emph{subfunctor}) of \(\mathbf{G}\) if \(\mathbf{F}(S)\) is a subset of
\(\mathbf{G}(S)\) for each \(S\).

\emph{In \(\widehat{\mathcal{C}}\) ``arbitrary'' projective limits exist
and are computed by:}%
{\renewcommand{\thefootnote}{(6)}\nde{Likewise, ``arbitrary'' inductive
limits exist and are computed ``argument by argument'', \ie
\((\varinjlim_{i}\mathbf{F}_{i})(S)=\varinjlim_{i}\mathbf{F}_{i}(S)\);
but in general the functor \(\mathbf{h}\) does not commute with inductive
limits.}}
\[
\bigl(\varprojlim_{i}\mathbf{F}_{i}\bigr)(S)
=\varprojlim_{i}\mathbf{F}_{i}(S).
\]
In particular, fibered products are defined
by:%
{\renewcommand{\thefootnote}{(7)}\nde{In particular, the kernel of a pair
of morphisms \(u,v\colon\mathbf{F}\rightrightarrows\mathbf{G}\) is the
subfunctor \(\Ker(u,v)\) of \(\mathbf{F}\) defined by
\(\Ker(u,v)(S)=\{x\in\mathbf{F}(S)\mid u(x)=v(x)\}\).}}
\[
\bigl(\mathbf{F}\times_{\mathbf{G}}\mathbf{F}'\bigr)(S)
=\mathbf{F}(S)\times_{\mathbf{G}(S)}\mathbf{F}'(S).
\]

We shall choose as final object of \(\widehat{\mathcal{C}}\) the functor
\(\mathbf{e}\) such that \(\mathbf{e}(S)=\{\varnothing\}\)\/.%
{\renewcommand{\thefootnote}{(8)}\nde{\(\{\varnothing\}\) (the set of
subsets of the empty set) denotes the one-element set.}}
Every \(\mathbf{F}\in\Ob(\widehat{\mathcal{C}})\) possesses a unique
morphism into \(\mathbf{e}\), and one sets
\[
\mathbf{F}\times\mathbf{F}'=\mathbf{F}\times_{\mathbf{e}}\mathbf{F}'.
\]

\emph{The functor \(\mathbf{h}\) commutes with projective limits; in
particular, for \(X\times X'\) to exist
(\(X,X'\in\Ob(\mathcal{C})\)), \resp for \(\mathcal{C}\) to admit a final
object \(e\), it is necessary and sufficient that
\(\mathbf{h}_{X}\times\mathbf{h}_{X'}\) be representable, \resp that
\(\mathbf{e}\) be representable, and one has}
\[
\mathbf{h}_{X}\times\mathbf{h}_{X'}\simeq\mathbf{h}_{X\times X'}
\qquad\text{and}\qquad
\mathbf{h}_{e}\simeq\mathbf{e}.
\]

\emph{A monomorphism of \(\widehat{\mathcal{C}}\) is nothing other than a
morphism \(\mathbf{F}\to\mathbf{G}\) such that, for
\(S\in\Ob(\mathcal{C})\), the corresponding map of sets
\(\mathbf{F}(S)\to\mathbf{G}(S)\) is injective.}%
{\renewcommand{\thefootnote}{(9)}\nde{If
\(\mathbf{F}(S)\to\mathbf{G}(S)\) is injective for every \(S\), it is
clear that \(\mathbf{F}\to\mathbf{G}\) is a monomorphism; the converse
is seen by considering the diagram
\(\mathbf{F}\times_{\mathbf{G}}\mathbf{F}\rightrightarrows\mathbf{F}\to\mathbf{G}\).
One thus obtains that: ``\(\mathbf{F}\to\mathbf{G}\) is a monomorphism
if and only if the diagonal morphism
\(\mathbf{F}\to\mathbf{F}\times_{\mathbf{G}}\mathbf{F}\) is an
isomorphism'' (\cf EGA~I, 5.3.8). Likewise, it is clear that if
\(\mathbf{F}(S)\to\mathbf{G}(S)\) is surjective for every \(S\), then
\(\mathbf{F}\to\mathbf{G}\) is an \emph{epimorphism}, and the converse
is seen by considering the amalgamated sum
\(\mathbf{G}\amalg_{\mathbf{F}}\mathbf{G}\), \cf the proof of
Lemma~4.4.4 in Exp.~IV.}}

\emph{The functor \(\Gamma\).}
For every \(\mathbf{F}\in\Ob(\widehat{\mathcal{C}})\) one sets
\[
\Gamma(\mathbf{F})=\Hom(\mathbf{e},\mathbf{F});
\]
an element of \(\Gamma(\mathbf{F})\) is thus a family
\((\gamma_{S})_{S\in\Ob(\mathcal{C})}\),
\(\gamma_{S}\in\mathbf{F}(S)\), such that for every arrow
\(f\colon S'\to S''\) of \(\mathcal{C}\), one has
\(\mathbf{F}(f)(\gamma_{S''})=\gamma_{S'}\).

One sets \(\Gamma(X)=\Gamma(\mathbf{h}_{X})\) for
\(X\in\Ob(\mathcal{C})\). If \(\mathcal{C}\) has a final object \(e\), one
thus has an isomorphism \(\Gamma(X)\simeq\Hom(e,X)\).

\numpara{1.3}
Let \(S\in\Ob(\mathcal{C})\). One denotes by \(\mathcal{C}/S\) the
category of objects of \(\mathcal{C}\) over \(S\), \ie the category whose
objects are the arrows \(f\colon T\to S\) of \(\mathcal{C}\), the set
\(\Hom(f,f')\) being the subset of \(\Hom(T,T')\) formed of the \(u\) such
that \(f=f'\circ u\). If \(\mathcal{C}\) possesses a final object \(e\),
then \(\mathcal{C}/e\) is isomorphic to \(\mathcal{C}\). The category
\(\mathcal{C}/S\) possesses a final object: the identity arrow
\(S\to S\).

If \(f\colon T\to S\) is an object of \(\mathcal{C}/S\), then one can
form the category \((\mathcal{C}/S)/f\), which one denotes, by abuse of
language, \((\mathcal{C}/S)/T\), and one has a canonical isomorphism
\[
\mathcal{C}/T\simeq(\mathcal{C}/S)/T.
\]
This construction also applies to the category \(\widehat{\mathcal{C}}\);
one defines in particular the category
\(\widehat{\mathcal{C}}/\mathbf{h}_{S}\). On the other hand, one can
form the category \(\widehat{\mathcal{C}/S}\).

If \(f\colon T\to S\) is an object of \(\mathcal{C}/S\), then
\(\Gamma(f)\) is identified with the set \(\Gamma(T/S)\) of sections of
\(T\) over \(S\), that is, of the arrows \(S\to T\) which are right
inverses of \(f\). Let us observe that
\(\mathbf{h}_{f}\colon\mathbf{h}_{T}\to\mathbf{h}_{S}\) is then an
object of \(\widehat{\mathcal{C}}/\mathbf{h}_{S}\) and that one has:
\[
\Gamma(\mathbf{h}_{f})\simeq\Gamma(\mathbf{h}_{T}/\mathbf{h}_{S})
\simeq\Gamma(T/S)\simeq\Gamma(f).
\]

\numpara{1.4}
One now proposes to define an equivalence of the categories
\(\widehat{\mathcal{C}/S}\) and \(\widehat{\mathcal{C}}/\mathbf{h}_{S}\),
that is, to prove that ``to give a functor on the category of objects
of \(\mathcal{C}\) over \(S\), is ``the same thing'' as to give a functor
on \(\mathcal{C}\) equipped with a morphism into
\(\mathbf{h}_{S}\)''.

\begin{enumeratei}
\item \emph{Construction of
\(\alpha_{S}\colon\widehat{\mathcal{C}}/\mathbf{h}_{S}
\to\widehat{\mathcal{C}/S}\).}

Let first \(\mathbf{H}\colon\mathbf{F}\to\mathbf{h}_{S}\) be an object of
\(\widehat{\mathcal{C}}/\mathbf{h}_{S}\). One must define a functor
\(\alpha_{S}(\mathbf{H})\) on \(\mathcal{C}/S\). Let therefore first
\(f\colon T\to S\) be an object of \(\mathcal{C}/S\); let us define
\(\alpha_{S}(\mathbf{H})(f)\) to be the inverse image of
\(f\in\mathbf{h}_{S}(T)\) under the map
\(\mathbf{H}(T)\colon\mathbf{F}(T)\to\mathbf{h}_{S}(T)\).%
{\renewcommand{\thefootnote}{(10)}\nde{For example, if
\(\mathbf{F}=\mathbf{h}_{X}\), then \(\mathbf{H}\) corresponds to a
morphism \(h\colon X\to S\) and
\(\mathbf{H}(T)\colon\Hom_{\mathcal{C}}(T,X)\to\Hom_{\mathcal{C}}(T,S)\)
is the map \(g\mto h\circ g\), whence
\(\alpha_{S}(\mathbf{h}_{X})=\Hom_{\mathcal{C}/S}(-,X)\).}}

Let next \(u\colon f\to f'\) be an arrow of \(\mathcal{C}/S\); then
\(\mathbf{F}(u)\colon\mathbf{F}(T')\to\mathbf{F}(T)\) induces a map from
\(\alpha_{S}(\mathbf{H})(f')\) to \(\alpha_{S}(\mathbf{H})(f)\), which one
denotes \(\alpha_{S}(\mathbf{H})(u)\). One checks at once that the maps
\[
f\mto\alpha_{S}(\mathbf{H})(f)
\qquad\text{and}\qquad
u\mto\alpha_{S}(\mathbf{H})(u)
\]
do define a functor on \(\mathcal{C}/S\), hence an object
\(\alpha_{S}(\mathbf{H})\) of \(\widehat{\mathcal{C}/S}\).

Let finally \(\mathbf{H}\colon\mathbf{F}\to\mathbf{h}_{S}\) and
\(\mathbf{H}'\colon\mathbf{F}'\to\mathbf{h}_{S}\) be two objects of
\(\widehat{\mathcal{C}}/\mathbf{h}_{S}\) and
\(\mathbf{U}\colon\mathbf{H}\to\mathbf{H}'\) a morphism of
\(\widehat{\mathcal{C}}/\mathbf{h}_{S}\):
\[
\xymatrix{
\mathbf{F}\ar[rr]^-{\mathbf{U}}\ar[rd]_-{\mathbf{H}}
& & \mathbf{F}'\ar[ld]^-{\mathbf{H}'}\\
& \mathbf{h}_{S} &
}
\]
Then for every \(f\colon T\to S\), the map
\(\mathbf{U}(T)\colon\mathbf{F}(T)\to\mathbf{F}'(T)\) induces a map
\[
\alpha_{S}(\mathbf{U})(f)\colon
\alpha_{S}(\mathbf{H})(f)\to\alpha_{S}(\mathbf{H}')(f),
\]
which defines a morphism of functors
\[
\alpha_{S}(\mathbf{U})\colon
\alpha_{S}(\mathbf{H})\to\alpha_{S}(\mathbf{H}').
\]
One checks easily that the maps
\[
\mathbf{H}\mto\alpha_{S}(\mathbf{H})
\qquad\text{and}\qquad
\mathbf{U}\mto\alpha_{S}(\mathbf{U})
\]
do define a functor
\(\alpha_{S}\colon\widehat{\mathcal{C}}/\mathbf{h}_{S}
\to\widehat{\mathcal{C}/S}\).

\item
\begin{proposition}{1.4.1}
\label{I.1.4.1}
The functor
\(\alpha_{S}\colon\widehat{\mathcal{C}}/\mathbf{h}_{S}
\to\widehat{\mathcal{C}/S}\)
is an equivalence of categories.
\end{proposition}
\end{enumeratei}

Let us merely indicate the principle of the construction of a
quasi-inverse functor
\(\beta_{S}\colon\widehat{\mathcal{C}/S}
\to\widehat{\mathcal{C}}/\mathbf{h}_{S}\).
Let \(\mathbf{G}\) be a functor on \(\mathcal{C}/S\); for every object
\(T\) of \(\mathcal{C}\), one sets
\[
\beta_{S}(\mathbf{G})(T)
=\text{sum of the sets \(\mathbf{G}(f)\)
for \(f\in\Hom(T,S)=\mathbf{h}_{S}(T)\)},
\]
which defines a functor \(\beta_{S}(\mathbf{G})\) on \(\mathcal{C}\),
equipped with an obvious projection onto \(\mathbf{h}_{S}\).

\numpara{1.5}
\emph{The equivalence \(\alpha_{S}\) commutes with the functors
\(\Gamma\).}
In other words, if \(\mathbf{H}\colon\mathbf{F}\to\mathbf{h}_{S}\) is an
object of \(\widehat{\mathcal{C}}/\mathbf{h}_{S}\) and
\(\alpha_{S}(\mathbf{H})\) the corresponding object of
\(\widehat{\mathcal{C}/S}\), one has
\[
\Gamma(\alpha_{S}(\mathbf{H}))\simeq\Gamma(\mathbf{H})
\simeq\Gamma(\mathbf{F}/\mathbf{h}_{S}).
\]
\emph{The equivalence \(\alpha_{S}\) commutes with the functors
\(\mathbf{h}\), \ie if \(f\colon T\to S\) is an object of
\(\mathcal{C}/S\),
\(\mathbf{h}_{f}\colon\mathbf{h}_{T}\to\mathbf{h}_{S}\) is an object of
\(\widehat{\mathcal{C}}/\mathbf{h}_{S}\) whose transform by
\(\alpha_{S}\) is none other than \(\mathbf{h}_{\mathcal{C}/S}(f)\),
where}
\[
\mathbf{h}_{\mathcal{C}/S}\colon\mathcal{C}/S
\to\widehat{\mathcal{C}/S}
\]
is the canonical functor%
{\renewcommand{\thefootnote}{(11)}\nde{\Cf N.D.E.~(10).}}.
Consequently:

\begin{proposition}{1.5.1}
\label{I.1.5.1}
Let \(\mathbf{H}\colon\mathbf{F}\to\mathbf{h}_{S}\) be an object of
\(\widehat{\mathcal{C}}/\mathbf{h}_{S}\). For
\(\alpha_{S}(\mathbf{H})\colon(\mathcal{C}/S)^{\circ}\to\Ens\) to be
representable, it is necessary and sufficient that
\(\mathbf{F}\colon\mathcal{C}^{\circ}\to\Ens\) be representable; if
\(\mathbf{F}\simeq\mathbf{h}_{T}\), then \(\alpha_{S}(\mathbf{H})\) is
representable by the object \(T\to S\) of \(\mathcal{C}/S\).
\end{proposition}

\emph{The equivalence \(\alpha_{S}\) is transitive in \(S\): if
\(f\colon T\to S\) is an object of \(\mathcal{C}/S\), one has a
commutative diagram of equivalences}
\[
\xymatrix@C=2.6em@R=1.8em{
(\widehat{\mathcal{C}}/\mathbf{h}_{S})/\mathbf{h}_{T}
  \ar[rr]^-{\alpha_{S}/\mathbf{h}_{T}}
  \ar[rd]_-{\simeq}
&&
\widehat{(\mathcal{C}/S)}/\mathbf{h}_{T}
  \ar[rr]^-{\alpha_{f}}
&&
\widehat{(\mathcal{C}/S)}/T
  \ar[ld]^-{\simeq}
\\
&
\widehat{\mathcal{C}}/\mathbf{h}_{T}
  \ar[rr]_-{\alpha_{T}}
&&
\widehat{\mathcal{C}/T}
&
},
\]
where \(\alpha_{S}/\mathbf{h}_{T}\) denotes (provisionally) the
\emph{restriction} (\cf 1.6) of the functor \(\alpha_{S}\) to the objects
over \(\mathbf{h}_{T}\).

\par\addvspace{0.55\baselineskip}
\noindent\textbf{1.6. Base change in a functor.}\ ---
For every \(S\in\Ob(\mathcal{C})\), one has a canonical functor
\[
i_{S}\colon\mathcal{C}/S\to\mathcal{C}
\]
defined by \(i_{S}(f)=T\) if \(f\) is the arrow \(T\to S\). If
\(f\colon T\to S\) is an object of \(\mathcal{C}/S\), one denotes
\(i_{T/S}=i_{f}\) the functor:
\[
i_{T/S}\colon(\mathcal{C}/S)/T\to\mathcal{C}/S,
\]
and one has the commutative diagram:
\[
\xymatrix@C=3.2em{
(\mathcal{C}/S)/T
  \ar[r]^-{i_{T/S}}
  \ar[rd]_-{\simeq}
&
\mathcal{C}/S
  \ar[r]^-{i_{S}}
&
\mathcal{C}
\\
&
\mathcal{C}/T
  \ar[ru]_-{i_{T}}
&
}.
\]
\emph{That is, identifying \((\mathcal{C}/S)/T\) with \(\mathcal{C}/T\) as
we shall do henceforth,}
\[
i_{S}\circ i_{T/S}=i_{T}.
\]
In the same way, if one identifies \(\mathcal{C}\) and \(\mathcal{C}/e\),
when \(\mathcal{C}\) possesses a final object \(e\), then
\(i_{S/e}\colon\mathcal{C}/S\to\mathcal{C}/e\) is identified with
\(i_{S}\).

For \(X\in\Ob(\mathcal{C})\) (\resp \(Y\in\Ob(\mathcal{C}/S)\)), let
\(p_{S}(X)\) (\resp \(p_{T/S}(Y)\)) be the object of \(\mathcal{C}/S\)
(\resp of \(\mathcal{C}/T\)) when it exists, defined by \(X\times S\)
(\resp \(Y\times_{S}T\)) equipped with its second projection:
\[
\vcenter{\xymatrix@R=1.5em{
X\times S\ar[d]_{p_{S}(X)}\\
S
}}
\qquad\text{resp.}\qquad
\vcenter{\xymatrix@R=1.5em{
Y\times_{S}T\ar[d]^{p_{T/S}(Y)}\\
T
}}.
\]

The (partially defined) functor \(p_{S}\) (\resp \(p_{T/S}\)) is called
the \emph{base change functor}. It is, by definition of the product
(\resp of the fibered product), the functor right adjoint to the functor
\(i_{S}\) (\resp \(i_{T/S}\))%
{\renewcommand{\thefootnote}{(12)}\nde{I.e., if \(g\colon U\to S\)
(\resp \(h\colon V\to T\)) is an object of \(\mathcal{C}/S\)
(\resp \(\mathcal{C}/T\)), then \(i_{S}(g)=U\) and \(i_{T/S}(h)\) is the
object \(f\circ h\colon V\to S\) of \(\mathcal{C}/S\), and one has
\(\Hom_{\mathcal{C}/S}(U,X\times S)\simeq\Hom_{\mathcal{C}}(U,X)\),
resp.\
\(\Hom_{\mathcal{C}/T}(V,X\times_{S}T)\simeq\Hom_{\mathcal{C}/S}(V,X)\).}}%
.
One likewise denotes
\[
p_{S}(X)=X_{S}
\qquad\text{and}\qquad
p_{T/S}(Y)=Y_{T}.
\]
The functor \(i_{S}\) defines a functor (\emph{restriction})
\[
i_{S}^{*}\colon\widehat{\mathcal{C}}\to\widehat{\mathcal{C}/S};
\]
one denotes \(\mathbf{F}_{S}=i_{S}^{*}(\mathbf{F})=\mathbf{F}\circ i_{S}\).
One evidently has
\[
i_{T/S}^{*}\circ i_{S}^{*}=i_{T}^{*},
\]
that is, for every functor \(\mathbf{F}\in\Ob(\widehat{\mathcal{C}})\),
\[
(\mathbf{F}_{S})_{T}=\mathbf{F}_{T}.
\]
The notation calls for a justification, which is as follows:

\begin{proposition}{1.6.1}
\label{I.1.6.1}
For the functor
\((\mathbf{h}_{X})_{S}\colon(\mathcal{C}/S)^{\circ}\to\Ens\)
to be representable, it is necessary and sufficient that the product
\(X\times S\) exist. One then has
\[
(\mathbf{h}_{X})_{S}\simeq\mathbf{h}_{X_{S}}.
\]
\end{proposition}

This shows that \(\mathbf{F}_{S}\) has two interpretations:
\emph{restriction} of the functor \(\mathbf{F}\) to \(\mathcal{C}/S\),
functor obtained by \emph{base change} \(\mathbf{e}\leftarrow S\). This
leads to the following notation:
\[
\xymatrix@R=1.8em@C=3em{
\mathbf{F} & \mathbf{F}_{S} & \mathbf{F}_{T}\\
\mathbf{e} & S\ar[l] & T\ar[l]
}
\]
which accounts well for the two preceding interpretations.

Let us observe that one has
\[
\Gamma(\mathbf{F}_{S})\simeq\Hom(\mathbf{h}_{S},\mathbf{F})\simeq\mathbf{F}(S),
\]
in particular
\[
\Gamma(X_{S})\simeq\Hom(S,X).
\]

\numpara{1.7.0}
{\renewcommand{\thefootnote}{(13)}\nde{The lemma below has been added
(\cf SGA~4, I.3.4); it is used in the proof of 1.7.1 and will be useful
several times in the sequel.}}%
Let \(\mathbf{E}\) be an object of \(\widehat{\mathcal{C}}\). Let us
consider the category \(\mathcal{C}_{/\mathbf{E}}\) of objects of
\(\mathcal{C}\) over \(\mathbf{E}\): its objects are the pairs
\((V,\rho)\) consisting of an object \(V\) of \(\mathcal{C}\) and of a
\(\widehat{\mathcal{C}}\)-morphism
\(\rho\colon\mathbf{h}_{V}\to\mathbf{E}\), \ie
\(\rho\in\mathbf{E}(V)\); a morphism from \((V,\rho)\) to
\((V',\rho')\) is the datum of a \(\mathcal{C}\)-morphism
\(f\colon V\to V'\) such that \(\rho'\circ f=\rho\)
(\ie \(\mathbf{E}(f)(\rho')=\rho\)).
